\subsection{Weak Lensing Signatures in Voids}
  \label{subsec:void_lensing}

  Because the effective framework modifies the inferred gravitational response
  in low-density regions, it also impacts weak gravitational lensing by cosmic
  voids.

  The model predicts subtle, environment-dependent modifications in the
  inferred lensing convergence profiles of deep voids.
  We emphasize that no global inconsistency with weak lensing surveys
  is implied.
  In standard-density environments, where $g_N \gtrsim a_\star$,
  the framework reduces to the unsaturated regime and reproduces
  conventional expectations.
  Potential deviations are restricted to the extreme low-acceleration
  outskirts of large voids, where the effective response approaches
  its saturation limit.

  In this regime, the effect would manifest not as a breakdown of
  lensing consistency, but as a mild shift in the outer convergence
  profile relative to purely kinematic mass inferences.
  Such deviations remain within current observational uncertainties,
  but provide a discriminant target for future high-precision
  surveys probing the deepest underdense regions~\cite{Cai2017VoidLensing}.

  Void lensing measurements therefore provide an independent probe of the
  framework, complementary to kinematic tests, and have been identified as
  sensitive diagnostics of gravitational behavior in underdense
  regions~\cite{Cai2017VoidLensing}.
